\documentclass[11pt]{article}
\usepackage[margin=0.8in]{geometry}
\usepackage[T1]{fontenc}
\usepackage{lmodern,amsmath,amssymb,amsthm,booktabs,array}
\usepackage[colorlinks=true,linkcolor=blue,urlcolor=blue,citecolor=blue]{hyperref}
\hypersetup{pdftitle={A complete irreducibility classification and three recurrence refutations},pdfauthor={Blake Stone},pdfsubject={A327315, A289001, A289035, A289239, and complete Lean proofs}}
\setlength{\parindent}{0pt}
\setlength{\parskip}{0.38em}
\setlength{\emergencystretch}{1em}
\newtheorem{theorem}{Theorem}[section]
\newtheorem{lemma}[theorem]{Lemma}
\newtheorem{proposition}[theorem]{Proposition}
\newtheorem{corollary}[theorem]{Corollary}
\newcommand{\N}{\mathbb N}
\newcommand{\Z}{\mathbb Z}
\newcommand{\Q}{\mathbb Q}
\newcommand{\R}{\mathbb R}
\begin{document}
\begin{center}
{\LARGE\bfseries A complete irreducibility classification\\[3pt]
and three recurrence refutations}\par
\vspace{0.5em}
{\large Blake Stone}\par
Proof developed with ChatGPT\par
November 27, 2026
\end{center}

\textbf{Abstract.}
We prove that the factorial-normalized numerator of the derivative of
order $n$ of $(X-2)/(X^2-X+1)$ is irreducible over $\Q$ exactly when
$n+1=2^a3^b$ for nonnegative integers $a,b$.  A sixth-root quadratic
field and a cyclotomic degree argument prove irreducibility in every
allowed degree.  Explicit rational row factors, sometimes reflected
under $X\mapsto1-X$, handle every excluded degree.
We also refute the proposed common length sequence for three binary
replacement systems.  Their literal lengths agree with the proposed
sequence through iteration 21 and first disagree at iteration 22,
where the lengths are $268184$ instead of $268182$.  We prove exact
count dynamics, a uniform discrepancy bound, and the successive-length
ratio limit $1+\sqrt2/2$.  This nonintegral algebraic growth constant
rules out every eventual monic integer linear relation, of any order.
Two Lean projects, with sixteen local proof modules, verify these
statements from the literal source definitions.  Established polynomial
families, methods, and related word results are credited.  Bounded
literature and repository searches located no earlier proof of the
exact assertions addressed here; this is not a certification of priority.

\begin{center}
\begin{tabular}{@{}l>{\raggedright\arraybackslash}p{4.5in}@{}}
\toprule
Source & Result and scope\\
\midrule
A327315/A069353 & One complete irreducibility classification for every derivative order.\\
A289001/A289004 & One length-conjecture refutation for $01\mapsto001$.\\
A289035 & One length-conjecture refutation for $01\mapsto010$.\\
A289239 & One length-conjecture refutation for $01\mapsto100$.\\
\bottomrule
\end{tabular}
\end{center}

The three word assertions share a single count theorem; their repeated
appearance in A289004 and A289001 is counted once.  The corrected growth
and recurrence obstruction are stronger conclusions within that same
family.  The separate prime question in A069353 and the other position
or fixed-point conjectures in the word entries are outside these claims.

\section{The complete irreducibility classification in A327315}
\label{sec:a327315}

\subsection{The literal source statement}

Let $P_n$ be the numerator of the $n$th derivative of
\[
 f(X)=\frac{X-2}{X^2-X+1},
\]
divided by $n!$, with the rational function reduced and its denominator
monic.  OEIS A327315~\cite{A327315} records these rows from $n=0$ and conjectures that
their irreducible indices are exactly A069353~\cite{A069353}, the numbers $2^a3^b-1$.
The entry attributes the conjecture to Clark Kimberling in November 2019.
Its title says degree $n$, but the executable definition and displayed
rows have degree $n+1$, beginning with $P_0=X-2$.  We follow those
unambiguous row and derivative indices.

\begin{theorem}[All derivative orders]
\label{thm:a327315}
For every integer $n\ge0$,
\[
 P_n\text{ is irreducible over }\mathbb Q
 \quad\Longleftrightarrow\quad
 n+1=2^a3^b\text{ for some }a,b\ge0.
\]
The same criterion holds for the unnormalized reduced numerator $n!P_n$.
\end{theorem}

There is one classification here.  The stronger irreducibility over a
quadratic field, the explicit factors in the excluded degrees, and the
source's degree correction are not counted as additional conjectures.

\subsection{Recurrence and exact reduced numerator}

Put $A=1-2X$ and $D=1-X+X^2$, and define
\[
 S_0=-1,\qquad S_1=X-2,\qquad
 S_{N+2}=A S_{N+1}-D S_N.
\]
The top-coefficient recurrence has initial values $-1,1$ and satisfies
$c_{N+2}=-2c_{N+1}-c_N$.  Consequently
\[
 \deg S_N=N,\qquad \operatorname{lc}(S_N)=(-1)^{N+1}.
\]
Differentiating the recurrence and using two-step induction gives
\[
 S_N'=-N S_{N-1}\quad(N\ge1).
\]
Combining this with the recurrence yields, for all $N\ge0$,
\begin{equation}
\label{eq:a327315-differential}
 D S_N'+N A S_N=N S_{N+1}.
\end{equation}
Since $D'=-A$, the quotient rule and
\eqref{eq:a327315-differential} prove inductively that
\begin{equation}
\label{eq:a327315-derivative}
 f^{(n)}(X)=\frac{n!S_{n+1}(X)}{D(X)^{n+1}}.
\end{equation}

The identities
\[
 4D-A^2=3,\qquad D-(X+1)(X-2)=3
\]
give Bezout certificates over $\mathbb Q$ for $(A,D)$ and $(S_1,D)$.
The polynomial $S_0=-1$ is a unit.  Modulo $D$, the recurrence says
$S_{N+2}\equiv A S_{N+1}$, so induction proves
$\gcd(S_N,D)=1$ for every $N$.  Thus the fraction in
\eqref{eq:a327315-derivative} has no polynomial cancellation.  Its
denominator is monic, and $n!$ is a nonzero rational unit.  Therefore its
canonical numerator is exactly $n!S_{n+1}$, and
\begin{equation}
\label{eq:a327315-source}
 P_n=S_{n+1}.
\end{equation}
This identifies the recurrence with the original reduced-numerator
definition, including $n=0$.

\subsection{The quadratic field and fractional-linear form}

Let $K=\mathbb Q(u)$ with $u^2-u+1=0$, and set $v=1-u$.
The polynomial $Y^2-Y+1=\Phi_6(Y)$ is irreducible over $\mathbb Q$, so
$[K:\mathbb Q]=2$.  Direct algebra gives
\[
 u+v=uv=1,\quad u^2=-v,\quad v^2=-u,\quad
 u^3=v^3=-1,\quad u^6=v^6=1.
\]
Also $u\ne v$, and neither $u$ nor $v$ is zero or one.
With $\alpha=u-X$ and $\beta=v-X$, we have
$\alpha+\beta=A$ and $\alpha\beta=D$.  The initial values and
recurrence therefore prove
\begin{equation}
\label{eq:a327315-closed}
 S_N=-v\alpha^N-u\beta^N
 =(-1)^{N+1}v\big((X-u)^N-v(X-v)^N\big).
\end{equation}

We use the following standard fractional-linear observation.  Over any
field, if $a\ne b$, $c\ne1$, and $Y^N-c$ is irreducible, then
\[
 (X-a)^N-c(X-b)^N
\]
is irreducible of degree $N$.  Indeed, if $x$ is a root of an irreducible
factor, then $x\ne b$ and
$y=(x-a)/(x-b)$ satisfies $y^N=c$.  Thus the extension containing $x$
has degree at least $N$, while the displayed polynomial has degree $N$.
Every irreducible factor must have full degree.  The Lean helper proves
this argument directly with an adjoint-root field.

\subsection{Irreducibility in every allowed degree}

\begin{lemma}
\label{lem:a327315-cyclotomic}
For $N=2^a3^b$, with $a,b\ge0$,
\[
 Y^{2N}-Y^N+1=\Phi_{6N}(Y).
\]
\end{lemma}
\begin{proof}
The initial identity is $\Phi_6(Y)=Y^2-Y+1$.  If a prime $p$ divides
$r$, the classical cyclotomic expansion formula is
$\Phi_r(Y^p)=\Phi_{rp}(Y)$.  Starting at $r=6$, apply this formula
$a$ times with $p=2$ and $b$ times with $p=3$.  Both required divisibilities
hold at every step, and the left side becomes the stated trinomial.
\end{proof}

\begin{lemma}
\label{lem:a327315-binomial}
If $N=2^a3^b$, then $Y^N-v$ is irreducible over $K$.
\end{lemma}
\begin{proof}
Choose an irreducible factor $q\in K[Y]$, of degree $d$, and let $y$ be
its root in $E=K[Y]/(q)$.  Then $y^N=v$, and hence
\[
 y^{2N}-y^N+1=v^2-v+1=0.
\]
By Lemma~\ref{lem:a327315-cyclotomic}, this is the irreducible polynomial
$\Phi_{6N}$ over $\mathbb Q$, of degree $2N$.  It is therefore the
minimal polynomial of $y$ over $\mathbb Q$.  The tower law gives
\[
 2N\le[E:\mathbb Q]=[K:\mathbb Q][E:K]=2d.
\]
Since $q$ divides a degree-$N$ polynomial, $d\le N$.  Thus $d=N$, proving
the assertion.
\end{proof}

Equation~\eqref{eq:a327315-closed} and the fractional-linear observation
now prove that $S_N$ is irreducible over $K$ for every $N=2^a3^b$.
It follows that $S_N$ is irreducible over $\mathbb Q$ as well: a proper
rational factorization would remain a proper factorization over $K$.

\subsection{Explicit proper factors in every excluded degree}

Let $N>0$ have a prime divisor $p\notin\{2,3\}$, and write $N=pm$.
Then $p\equiv1$ or $5\pmod6$, and $1\le m<N$.
For any $c$ satisfying $c^p=v$, the difference-of-powers identity gives
\[
 (X-u)^m-c(X-v)^m
 \ \mid\ (X-u)^{pm}-v(X-v)^{pm}
 \qquad\text{in }K[X].
\]
If $p\equiv1\pmod6$, choose $c=v$.  Equation~\eqref{eq:a327315-closed}
then gives $S_m\mid S_N$ over $K$.

If $p\equiv5\pmod6$, choose $c=u$, since $u^5=v$.  Substituting $1-X$
in the closed form gives the exact reflected identity
\[
 S_m(1-X)=-u\big((X-u)^m-u(X-v)^m\big).
\]
Consequently $S_m(1-X)\mid S_N$ over $K$.

Both proposed divisors belong to $\mathbb Q[X]$.  Divisibility descends
to $\mathbb Q[X]$ by polynomial division and injectivity of the
coefficient embedding.  Each divisor has degree $m$, strictly between
zero and $N$, so $S_N$ is reducible.

Finally, a positive integer with no prime divisor other than 2 and 3
has the form $2^a3^b$, by induction on its prime factors.  This proves
the converse for every positive degree $N$.  Taking $N=n+1$ in
\eqref{eq:a327315-source} proves Theorem~\ref{thm:a327315}, including
the initial linear row.  Multiplication by the nonzero constant $n!$
does not change irreducibility.

\subsection{Prior literature and formal scope}

The source polynomial construction has an exact antecedent in
L.~Remete's \emph{A generalization of simplest number fields and their
integral basis}, arXiv:2111.08341v1 (2021)~\cite{Remete}.  His Section 2 defines
\[
 f_m^{(N)}(Y)=\sum_{j=0}^N\binom NjY^j g(N-j),
\]
where $g$ is six-periodic with values
$(1,-m,-m-1,-1,m,m+1)$.  Comparison with the definitions gives
$S_N(X)=(-1)^{N+1}f_1^{(N)}(X-1)$.
Indeed, at $m=1$, $g(k)=u^{k+1}+v^{k+1}$, so the binomial theorem gives
$f_1^{(N)}(Y)=u(Y+u)^N+v(Y+v)^N$, which agrees with the closed form
after the stated shift and sign.
The family itself, its derivative structure, and fractional-linear
root symmetries are therefore not new constructions here.
Remete's Corollary 10 requires a prime $p\ne3$ with
$v_p(m^2+m+1)=1$.  At $m=1$ that condition fails, and the inspected
irreducibility subsection does not state the complete degree criterion.

Related classical methods occur in C.~Cass, J.~Cullinan,
A.~Rasmussen and D.~Trifunovski, \emph{Specializations of Generalized
Rikuna Polynomials}, International Journal of Number Theory 10 (2014),
585--600~\cite{Rikuna}, and J.~S.~Calcut, \emph{The Tangent Analogues of the Chebyshev
Polynomials} (2008)~\cite{Calcut}.  Their relevant definitions and irreducibility or
angle-factorization sections were checked.  The present proof uses
standard cyclotomic irreducibility and fractional-linear substitution;
it does not claim those techniques as new.

The checked A327315 export still labels the exact classification a
conjecture.  The primary-reference review, exact and semantic web
queries, trureturing catalog, and targeted global and ProveIt GitHub
code searches located no earlier proof of that exact classification.
This is a bounded review, not a guarantee of historical priority.
The companion source-check report preserves the immutable OEIS exports,
their hashes, the index convention, and the scope of the prior-art search.

The formal development has nine local Lean modules.  It proves the
recurrence identities, exact degree, quadratic-field relations,
cyclotomic identity, both divisibility branches, the all-degree
criterion, and the exact canonical numerator bridge.  It uses the pinned Lean and mathlib versions
listed in Section~\ref{sec:verification}.
The rational-function derivative and fractional-linear helper are
reused unchanged from earlier packages.  The theorem dependencies are
the standard \texttt{propext}, \texttt{Classical.choice}, and
\texttt{Quot.sound}; no admitted proof or custom axiom is used.

\section{A shared false length recurrence for three binary scanners}
\label{sec:a289004}

The OEIS entries A289001, A289035, and A289239 give three replacement rules
starting with $w_0=00$ and conjecture that their iterate lengths are the
independently defined sequence A289004
\cite{A289001,A289004,A289035,A289239}.
We refute all three mapping assertions and prove a structural obstruction:
their common actual length sequence satisfies no eventual monic
constant-coefficient integer recurrence, of any order.
The assertion repeated in A289004 and A289001 is counted once.

\subsection{The literal sources and the first disagreement}

Every rule has $00\mapsto0010$ and $10\mapsto010$; the remaining image is
as follows.
\begin{center}
\begin{tabular}{lll}
\hline
Source word & Image of $01$ & Formal variant\\
\hline
A289001 & $001$ & \texttt{left}\\
A289035 & $010$ & \texttt{middle}\\
A289239 & $100$ & \texttt{right}\\
\hline
\end{tabular}
\end{center}
The scanner consumes nonoverlapping matches from the original input, copies
an unmatched leading letter, and retains a final singleton. Replacements
are not rescanned during the same pass. A289035 explicitly specifies the
retained-final-letter convention.

The proposed comparison sequence has
\[
 (a_0,\ldots,a_9)=(2,4,7,11,18,31,52,89,151,257)
\]
and, for $n\ge10$,
\begin{equation}
 a_n=2a_{n-1}-a_{n-2}+2a_{n-4}-a_{n-5}
       +a_{n-7}-a_{n-8}-a_{n-10}.
 \label{eq:a289004:claimed}
\end{equation}
The source prose accidentally repeats index 9 for the last two initial
values. Its displayed terms and Mathematica program unambiguously put 151
at index 8 and 257 at index 9; these are the values used in the formal
independent definition.

\paragraph{The source refutations.}
The three scanners have the same length $L_n$ at every iteration. For
every $n<22$, $L_n=a_n$, while
\[
 L_{22}=268184\ne268182=a_{22}.
\]
Moreover, recurrence~\eqref{eq:a289004:claimed}, when applied to the actual
lengths, fails after every finite cutoff.

The first differing values are:
\begin{center}
\begin{tabular}{r|r|r}
$n$ & actual $L_n$ & A289004 $a_n$\\
\hline
19 & 53908 & 53908\\
20 & 92027 & 92027\\
21 & 157099 & 157099\\
22 & 268184 & 268182\\
23 & 457819 & 457812
\end{tabular}
\end{center}
This is a failure of the proposed recurrence beyond a long common initial
segment, rather than a change in the seed or the starting index.

\subsection{The exact all-iteration count equations}

Call a finite word good if every 1 is immediately followed by a 0.
Equivalently, it is a concatenation of the blocks $0$ and $10$.
Every nonempty good word ends in 0 and has no factor $11$.
The seed is good, and each of the three scanners preserves this property.
For the left variant, the only output block ending in 1 is the image
$001$ of $01$. In a good input, such a pair must be followed by a 0,
and the output of the remaining suffix starts in 0. This checks the only
potential boundary problem; all other relevant output blocks end in 0.
The formal preservation argument treats every good finite input word.

Let $O_n$ denote the number of ones in the $n$th word and put
$e_n=L_n\bmod2$. A good word of length $L_n$ is scanned into exactly
$\lfloor L_n/2\rfloor$ pairs, followed by $e_n$ final zeros.
Every consumed pair produces one 1. A $00$ pair has output length 4,
and a mixed pair has output length 3. Since every input 1 belongs to
a mixed pair, this gives the exact equations
\begin{equation}
 O_{n+1}=\left\lfloor\frac{L_n}{2}\right\rfloor,
 \qquad
 L_{n+1}=2L_n-O_n-e_n,
 \qquad (L_0,O_0)=(2,0).
 \label{eq:a289004:counts}
\end{equation}
The right-hand sides do not depend on the variant. Induction therefore
proves the common length and common ones-count assertions. Eliminating
$O_n$ yields a corrected scalar recurrence:
\begin{equation}
 L_{n+2}+\left\lfloor\frac{L_n}{2}\right\rfloor
       +(L_{n+1}\bmod2)=2L_{n+1},
 \qquad L_0=2,\quad L_1=4.
 \label{eq:a289004:corrected}
\end{equation}
Evaluating this proved count state and the independent recurrence
\eqref{eq:a289004:claimed} establishes both the first discrepancy and
agreement at every earlier index. The Lean proof does not presume
\eqref{eq:a289004:counts}: it derives it from the literal lists.

\subsection{Uniform discrepancy and growth}

Set
\[
 \mu=1-\frac{\sqrt2}{2},\qquad
 \lambda=1+\frac{\sqrt2}{2},\qquad
 S_n=L_n-(2+\sqrt2)O_n.
\]
From \eqref{eq:a289004:counts} and
$2O_{n+1}+e_n=L_n$, direct substitution gives
\begin{equation}
 S_{n+1}=\mu S_n+(1-\mu)e_n.
 \label{eq:a289004:error}
\end{equation}
Since $S_0=2$, $0<\mu<1$, and $0\le e_n\le1$, induction proves
\begin{equation}
 0\le S_n\le2\qquad(n\ge0).
 \label{eq:a289004:bound}
\end{equation}
Every good word has at most half its letters equal to 1. Combining
$2O_n\le L_n$ with \eqref{eq:a289004:counts} gives
\[
 L_{n+1}\ge L_n+\left\lfloor\frac{L_n}{2}\right\rfloor,
 \qquad L_n\ge n+2.
\]
Thus $L_n\to\infty$. A second rearrangement of the exact count equations is
\begin{equation}
 \frac{L_{n+1}}{L_n}
 =\lambda+\mu\frac{S_n}{L_n}-\frac{e_n}{L_n}.
 \label{eq:a289004:ratio}
\end{equation}
The two error quotients tend to zero by \eqref{eq:a289004:bound}, so
\begin{equation}
 \lim_{n\to\infty}\frac{L_{n+1}}{L_n}
 =1+\frac{\sqrt2}{2}.
 \label{eq:a289004:limit}
\end{equation}
This argument proves existence of the limit. For every fixed
$j\ge0$, multiplying $j$ successive ratios proves
$L_{n+j}/L_n\to\lambda^j$.

\subsection{The obstruction to every eventual integer recurrence}

\paragraph{Theorem.}
For every monic polynomial $p(X)=\sum_i p_iX^i\in\mathbb Z[X]$
and every cutoff $N$, there is an $n\ge N$ such that
\[
 \sum_i p_iL_{n+i}\ne0.
\]
In particular, no recurrence
$L_{n+d}=\sum_{i<d}c_iL_{n+i}$ with fixed $c_i\in\mathbb Z$
can hold eventually.

\paragraph{Proof.}
If such a relation held for every sufficiently large $n$, division by
$L_n$ and the fixed-shift limits would give
\[
 p(\lambda)=\sum_i p_i\lambda^i=0.
\]
Monicity would make $\lambda$ an algebraic integer. Closure under
subtraction and multiplication would then make
\[
 (\lambda-1)^2=\frac12
\]
an algebraic integer. The rational-root theorem says that a rational
algebraic integer is an integer, contradicting $1/2\notin\mathbb Z$.
\hfill$\square$

The characteristic polynomial of the precise source recurrence is
\[
 p_0(X)=X^{10}-2X^9+X^8-2X^6+X^5-X^3+X^2+1.
\]
It is monic, so the theorem proves failures of the source recurrence
arbitrarily far out, not merely at iteration 22. The coefficient scope
is exactly monic integer recurrences; arbitrary rational or real
coefficients are not covered by the stated result.

\subsection{Prior results and formal endpoints}

Dekking and Keane proved that the infinite fixed word A289001 is the
Pell word \cite[Theorem~4]{DekkingKeane}. The related position results
are already known. Dekking also proved a floor recurrence for the
final-letter-removed iteration in A293077 \cite{A293077}; the related
density discussion is recorded in A289037 \cite{A289037}.
Those are prior results, not additional new assertions of this work.
The removed-letter lengths differ from the literal retained-singleton
lengths already at iteration 2, where they are 6 rather than 7.

The current source entries, both web-search engines, the trureturing
catalogue and its relevant replacement report, the existing repository,
and global/ProveIt GitHub code searches were checked. No matching earlier
refutation of these retained-singleton length assertions was found in
that bounded screen. Unindexed or differently phrased predecessors
are not excluded.

The seven-module Lean package proves the entire source chain.
\texttt{word\_counts} is the all-index literal scanner bridge;
\texttt{first\_failure} certifies the least failing index;
\texttt{three\_source\_conjectures\_false} states the three mapping
refutations; and \texttt{literal\_corrected\_recurrence} states
\eqref{eq:a289004:corrected}. The analysis endpoints are
\texttt{discrepancy\_bounds} and
\texttt{literal\_length\_ratio\_tendsto}. The stronger source-facing
obstruction is
\begin{center}
\texttt{literal\_no\_eventual\_monic\_relation}.
\end{center}
The exact order-ten specialization is
\begin{center}
\texttt{original\_recurrence\_fails\_after\_every\_cutoff}.
\end{center}
The formalization uses the pinned versions in
Section~\ref{sec:verification}. Its main theorems use
only the standard axioms \texttt{propext}, \texttt{Classical.choice}, and
\texttt{Quot.sound}; no extra axiom or native-decision shortcut is used.

\section{Formal artifacts, verification, and attribution}
\label{sec:verification}

The two standalone Lake projects pin Lean 4.19.0 and mathlib revision
\begin{center}
\path{c44e0c8ee63ca166450922a373c7409c5d26b00b}.
\end{center}
All sixteen local proof modules are explicit build roots.  Each project
is built from source in an author directory and a separate independent
directory.  These builds copy no local proof objects; they reuse only
the pinned external dependency cache.  The package records include
the complete build output, exact source hashes, source snapshots, and
the scope of the prior-publication searches.

For the derivative family, the final source-facing equivalence is
\begin{center}
\path{TwistedDerivative.source_irreducible_iff}.
\end{center}
The companion endpoint
\path{TwistedDerivative.source_numerator_irreducible_iff}
gives the same criterion for the unnormalized canonical numerator.
Thus the classification does not rely on an unproved bridge from a
recurrence to the source rational function.

For the word family, \path{A289004.three_source_conjectures_false}
refutes the three exact mapping claims.  The literal-list limit and
recurrence obstruction are
\path{A289004.literal_length_ratio_tendsto} and
\path{A289004.literal_no_eventual_monic_relation}.
All reported theorem dependencies are drawn from the standard axioms
\path{propext}, \path{Classical.choice}, and \path{Quot.sound}.
No local admitted theorem, custom axiom, unsafe declaration, or
native-evaluation proof shortcut is used.

The mathematical reviews check source definitions and quantifiers in
addition to successful compilation.  The two package-level independent
reviews are performed by contributors who authored none of the package
they are reviewing.  An additional manuscript review compares the
conventional proofs with the formal endpoints.  These are assistant
reviews and kernel checks, not claims of journal peer review or
editorial acceptance by the OEIS.

\subsection*{Correction to the preceding batch}

During this batch's expanded repository search, we found a September
19, 2026 report in VladimirReshetnikov/ProveIt that already proves
the entire A211417 positive-shift integrality family and a stronger
polynomial-divisor classification~\cite{PriorA211417}.  The preceding
batch initially counted that assertion as a new resolution.  Its Lean
proof remains unchanged and valid, but the appropriate contribution
label is an independent formalization of prior work.  The accompanying
Batch 10 manuscript and guide now make that correction and retain
their original reviews as historical records.  This attribution
correction is not counted among the current batch's four source
assertions.  The dated correction record gives the precise prior
proposition, original repository commit, and before-and-after hashes.

\begin{thebibliography}{99}
\raggedright

\bibitem{A327315}
The OEIS Foundation Inc., \emph{A327315}, entry attributed to Clark
Kimberling, November 1, 2019; revision 6, November 6, 2019.
\url{https://oeis.org/A327315}.
Official export checked November 27, 2026; immutable export commit
\texttt{571197984aeda7b8b680fe3a1f6a43f9898cd441}.

\bibitem{A069353}
The OEIS Foundation Inc., \emph{A069353: Numbers of form $2^i3^j-1$
with $i,j\ge0$}, revision 20, April 22, 2025.
\url{https://oeis.org/A069353}.
Official export checked November 27, 2026.

\bibitem{Remete}
L.~Remete, \emph{A generalization of simplest number fields and their
integral basis}, arXiv:2111.08341v1, November 16, 2021.
\url{https://arxiv.org/abs/2111.08341}.

\bibitem{Rikuna}
C.~Cass, J.~Cullinan, A.~Rasmussen and D.~Trifunovski,
\emph{Specializations of Generalized Rikuna Polynomials},
International Journal of Number Theory \textbf{10}(3) (2014), 585--600.
DOI: 10.1142/S1793042113501108.
Author-hosted preprint:
\url{https://homepages.math.uic.edu/~dtrifu2/pdf/rikuna_spec.pdf}.

\bibitem{Calcut}
J.~S.~Calcut, \emph{The Tangent Analogues of the Chebyshev Polynomials},
November 10, 2008.
\url{https://www2.oberlin.edu/faculty/jcalcut/arctan.pdf}.

\bibitem{A289001}
Clark Kimberling, OEIS A289001, fixed point of
$00\mapsto0010$, $01\mapsto001$, $10\mapsto010$.
Revision 37, 9 March 2022.
\url{https://oeis.org/A289001}.

\bibitem{A289004}
Clark Kimberling, OEIS A289004, the order-ten recurrence and
conjectured replacement lengths.
Revision 15, 7 April 2020.
\url{https://oeis.org/A289004}.

\bibitem{A289035}
Clark Kimberling, OEIS A289035, fixed point of
$00\mapsto0010$, $01\mapsto010$, $10\mapsto010$.
Revision 21, 22 March 2022.
\url{https://oeis.org/A289035}.

\bibitem{A289239}
Clark Kimberling, OEIS A289239, fixed point of
$00\mapsto0010$, $01\mapsto100$, $10\mapsto010$.
Revision 7, 7 April 2020.
\url{https://oeis.org/A289239}.

\bibitem{A293077}
Clark Kimberling, OEIS A293077, lengths for the final-letter-removed
mapping at A289035; proof of its floor recurrence by Michel Dekking,
20 March 2022. Revision 26, 23 March 2022.
\url{https://oeis.org/A293077}.

\bibitem{A289037}
Clark Kimberling, OEIS A289037, positions of ones in A289035;
comments by Michel Dekking, March 2022.
Revision 13, 22 March 2022.
\url{https://oeis.org/A289037}.

\bibitem{DekkingKeane}
Michel Dekking and Michael Keane,
\emph{Two-block substitutions and morphic words},
arXiv:2202.13548v2, 20 February 2023.
\url{https://arxiv.org/abs/2202.13548}.

\bibitem{PriorA211417}
\emph{Polynomial divisibility of height-one factorial ratios},
Part I, manuscript dated September 19, 2026, author line OpenAI ChatGPT,
in VladimirReshetnikov/ProveIt.  In particular, Proposition
\path{prop:positive}, Corollary \path{cor:classificationAB}, and the
subsection on unrestricted positive offsets.
\href{https://github.com/VladimirReshetnikov/ProveIt/blob/a3fe9660ee6c3ab92752594c2099c3cd214012c9/docs/reports/congruences-and-valuations/factorial-ratio-polynomial-divisibility/article.tex}{Pinned original September 19 source},
commit \path{a3fe9660ee6c3ab92752594c2099c3cd214012c9}.
\end{thebibliography}
\end{document}
